% ch-a6-01.tex —— 第 6 篇（Tony Feng, Michael Harris）篇首 + §1 + §2.1–2.6
% 底本 PDF 484–495 页 = 印刷页 471–482
\chapter{Langlands 对应中的导出结构}
\ysrc{471}
\begin{center}
{\small\itshape Tony Feng \quad Michael Harris}
\end{center}
\medskip
\noindent\textbf{摘要.} 本文综述近期出现的若干导出结构实例，它们都与 Langlands 对应相关联。案例研究包括\tn{底本此处作「Cases studies」，疑为「Case studies」之误。}导出伽罗瓦形变环、导出 Hecke 代数、导出 Hitchin 叠与导出特殊闭链。我们还强调若干公开问题，我们预期它们对未来的进展至关重要。
\footnote{2020 年数学主题分类：主分类 11F70、11F75、11F80、11F85、11F27。}
\footnote{第一作者受 NSF 基金 DMS-1902927、DMS-2302520、DMS-2441922 以及 Simons 基金会的一项基金资助。第二作者部分受 NSF 基金 DMS-2302208、DMS-2001369 资助。\textcopyright 2025 作者所有。}

\medskip
\noindent\textbf{目录.}
\begin{list}{}{\leftmargin=2.2em\labelwidth=2em\labelsep=0.4em\itemindent=-2.2em}
\item 1.\quad 动机与概览
\item 2.\quad 导出代数几何导引
\item 3.\quad 与 Langlands 纲领相关的导出模空间
\item 4.\quad 导出特殊闭链与高阶 theta 函数
\item 5.\quad 导出伽罗瓦形变环
\item 6.\quad 导出 Hecke 代数
\item 7.\quad 局部对称空间的上同调
\item 8.\quad Venkatesh 的动机猜想
\item 9.\quad 若干公开问题选
\item 附录 A.\quad 单纯交换环速成
\item 参考文献
\end{list}

\section{动机与概览}

\subsection{「Langlands 对应中的导出结构」指什么？}

「导出」这个形容词如今被用来修饰数学中的各种构造，含义常常大不相同。例如「导出函子」「导出范畴」「导出代数几何」都是我们将会遇到的实例。一般说来，「导出」一词指的是对数学构造的一种增强（enhancement），这种增强融入了同伦论。过去十年间，同伦论在 Langlands 对应上的应用在若干不同方向上迅速发展。下面概述本综述将会涉及的几个方面。

\ysrc{472}
\subsubsection{导出 Hecke 代数}

人们最早遇到的「导出函子」例子或许是 Tor 与 Ext 函子，它们通过对 $\otimes$ 与 $\Hom$ 作导出而构造。在经典 Langlands 对应中，Hecke 算子起着核心作用，这些算子构成作用在自守形式上的 Hecke 代数。Hecke 代数可以看作某种自同态空间；把自同态换成「导出自同态」（即 Ext 群），就得到导出 Hecke 代数的概念。这一概念出现在两个颇为不同的背景中：Ollivier--Schneider 关于 $p$-进 Langlands 对应的工作 [OS19]，以及 Venkatesh 等人关于算术群上同调的工作 [Ven19,GV18,PV21]。

\subsubsection{导出模空间}

上述意义的导出函子属于经典同调代数的范围：在那里，人们对阿贝尔范畴（例如环上的模范畴）上的函子作导出。非阿贝尔范畴上也有导出函子理论，它可追溯到 Quillen 的同伦代数 [Qui67]。把它应用于交换环范畴，就进入「导出代数几何」的领域。

在导出代数几何中，人们为代数几何的通常构造建立导出增强。例如，我们将讨论「导出概形」「导出叠」等。可以把导出概形设想为经典概形再加上一些额外的同伦信息。

最近人们意识到，这类导出增强为 Langlands 对应中的结构提供了自然的解释。两类不同的例子是：Galatius--Venkatesh 的导出伽罗瓦形变空间 [GV18]，它们被用来解释算术群上同调的结构；以及 Feng--Yun--Zhang 的导出 Hitchin 叠 [FYZ21]，它们被用来为与 Kudla 纲领相关的特殊闭链构造虚拟基本闭链。

\subsubsection{其他方向（本文不涉及）}

我们再提一下 Langlands 对应的若干其他方面：导出概念在其中起关键作用，尽管本文不会详细展开。

\begin{itemize}[leftmargin=1.6em]
\item 几何 Langlands 对应以曲线上局部系统的模叠为特征，特别是其上的凝聚层范畴。要得到该范畴的「正确」版本——也就是具有 Langlands 对偶解释的那个版本——一般需要以导出的方式来理解这个模叠。这类问题在本文集中 Sam Raskin 的文章 [Ras25] 中处理。
\item 导出交换代数在 Bhatt--Morrow--Scholze [BMS19] 与 Bhatt--Scholze [BS22] 近期关于整 $p$-进 Hodge 理论的工作中起重要的技术作用。
\end{itemize}

\subsection{为什么要导出？}

导出代数几何需要相当多的技术背景，因此自然会问：对于关心 Langlands 对应中较经典方面的研究者——例如自守形式与伽罗瓦表示之间的互反律——它为何有用、又如何有用？\footnote{我们将几何 Langlands 理论这类分支排除在外：在那类分支里，导出代数几何已经内置于核心问题的表述之中，因而其相关性显而易见。}

\ysrc{473}
这里有若干彼此独立的理由，本综述会涉及其中一部分。其一：局部对称空间的高阶上同调已成为 Langlands 对应的核心，因为上同调为自守形式空间提供自然的整结构，而引入 $p$-进方法正需要整结构。在许多经典情形中，相关的上同调群集中在单一次数，这使人可以在很大程度上忽略导出方面；但一般说来，它们分布在许多次数上，于是就必须使用复形、导出函子、导出范畴等工具。

到此为止的讨论只需要导出阿贝尔结构，这套理论自 20 世纪 70 年代以来已发展成熟并获广泛使用。但为了联系到伽罗瓦表示，我们想引用伽罗瓦表示的模空间，即所谓伽罗瓦形变环。要得到这些对象「相应地分布在许多次数上」的版本，就必须导出交换环的概念。这是导出代数几何的起点，需要更精细的同伦方法。

本文处理的第三条独立线索，与自守形式相关的计数几何有关。一个著名的例子是 Kudla 纲领，它力求在算术几何中实现 theta 函数的一种化身。更一般地，人们希望把 [BZSV24] 意义上的相对 Langlands 对偶在算术几何中实现出来。这意味着，我们寻求把自守形式实现为志村簇与 shtuka 模叠的 Chow 群或高阶上同调群中的闭链类。我们在本文中说明，这类闭链类应当来自导出几何：经典代数几何中的闭链一般性态很差，而它们的导出版本具有所希望的性质。

还有一些动机本文不涉及，包括 Xinwen Zhu [Zhu25] 与 Emerton--Gee--Hellman [EGH25] 的文章中所解释的 shtuka 模叠上同调的猜想性描述。在这些文章讨论的某些情形中，为了得到「正确」的答案，必须把导出结构纳入进来。

\subsection{本文的写法}

正如上述概述想要传达的，过去几年里，同伦论以 $\infty$-范畴、导出代数几何、稳定同伦论等形式迅速而异常广泛地涌入数论。

这套理论的技术基础令人生畏，文献中众多不同的进路又使之更加复杂。例如，文献分为模型范畴与 $\infty$-范畴两种语言（后者近来更流行）；即使在 $\infty$-范畴的框架内，也有不同的进路（拟范畴、DG-范畴、单纯范畴）。导出代数几何也有不同的风味，可以建立在交换微分分次代数、单纯交换环与 $E_\infty$-代数之上。这些问题很重要，但常常与应用这些工具于数论时所遇到的问题完全无关。

因此，鉴于同伦论背景方面已有大量文献，我们选择把本文写成一份「使用手册」，讲它如何应用于数论。我们甚至不打算对基础材料作详尽的技术处理。偶尔，我们会指出\ysrc{474}存在不同的进路，并评论它们在「直观」上的差异。多数时候，我们让事情保持一定程度的含糊与非正式，依靠类比与直觉，而非形式定义。

\subsection{本文的结构}

在第 2 节，我们对导出代数几何作一个精神性的导引。我们不提供任何技术细节——事实上，我们几乎没有作出任何数学上精确的陈述；我们只是试图对导出代数几何中出现的某些词语和短语，给出「如何思考它们」的指引。附录 A 对此作了平衡：它纯粹聚焦于导出代数的基本定义。

在第 3 节，我们讨论导出模空间的概念——正是它构成了我们最初使用导出代数几何的 raison d'être（存在理由）。我们试图给出一套启发式原则的一般模式，供实践中判断何时应当考虑导出模空间、以及如何构造它们。我们用在 Langlands 对应中具有相关性的若干例子来说明这些启发式原则。

在第 4 节，我们更详细地综述与 Kudla 纲领相关的导出模空间的出现与应用，重点放在 Feng--Yun--Zhang [FYZ24,FYZ21,FYZ23] 关于函数域上酉群的高阶 theta 级数的工作。

在第 5 节，我们以极其非正式的方式综述 Galatius--Venkatesh [GV18] 关于导出伽罗瓦形变环的工作。

在第 6 节，我们介绍局部导出 Hecke 代数，涉及 [Ven19] 中 $\ell\neq p$ 的结果，并简要提及 Schneider [Sch15a]、Ollivier 与 Ronchetti 关于 $\ell=p$ 的结果。

在第 7 节，我们讨论局部对称空间的上同调与整体导出 Hecke 代数；第 8 节解释 Venkatesh 的动机猜想。

最后，第 9 节提出若干处于导出代数几何与 Langlands 纲领交界处的公开问题，它们应当会有重要的推论。

\subsection{致谢}

我们感谢 Henri Darmon、Dennis Gaitsgory、Soren Galatius、Chandrashekhar Khare、Barry Mazur、Arpon Raksit、Victor Rotger、Akshay Venkatesh、Jonathan Wang、Zhiwei Yun 与 Wei Zhang 与我们讨论本材料，并感谢允许我们在此介绍的合作成果。我们感谢 Gurbir Dhillon、Justin Wu 与匿名审稿人对草稿的评论与修正。

附录 A 的初稿曾由第一作者在 2022 年春季的课程中向 MIT 研究生讲授试用。我们感谢参与者们的反馈。

\section{导出代数几何导引}

\subsection{什么是导出代数几何？}

代数几何研究由交换环搭建起来的几何。导出代数几何则研究以类似方式、由一些我们暂时非正式地称为「导出交换环」的对象搭建起来的几何。导出交换环是交换环的某种增强；它们支配着人们或许会称为「导出仿射概形」的东西，而后者构成导出概形的局部积木。因此，从代数几何到导出代数几何的过渡，主要就是把交换环换成导出交换环。

\ysrc{475}
\subsection{什么是导出交换环？}

一开始要提醒：对于「导出交换环」这一概念，有多个不同的进路。它们都相当技术化，因此我们在此不进入其中任何一个的形式定义，而是满足于解释直觉、类比，以及我们在实践中如何思考这些事情。

\subsubsection{单纯交换环}

本文大部分内容中，我们所谓的「导出交换环」将指单纯交换环（simplicial commutative ring）。「单纯」这个形容词可以看作「拓扑的」的同义词，所以一个合理的第一直觉是：把单纯交换环想成拓扑交换环。\footnote{我们避免直接使用拓扑交换环（这当然是定义良好的数学概念）的一个原因，是想排除「病态」的拓扑空间。因此对「单纯」更精细的口号或许是「拓扑的且性质良好的」。}

特别地，单纯交换环 $R$ 带有（阿贝尔的）同伦群 $\pi_0(R),\pi_1(R),\pi_2(R)$ 等。应当把它们想成 $R$ 底层的「拓扑空间」的同伦群；不仅如此，$R$ 上的交换环结构还使 $\pi_*(R)$ 具有分次交换环的结构。理解这个分次交换环，就近似于「理解」$R$——正如理解一个拓扑空间的通常拓扑同伦群，近似于「理解」这个空间。

单纯交换环的技术性定义相当繁复，因此被隔离到附录 A。本文正文写成这样：把「导出交换环」这一概念当作黑箱，读者也能读懂；第一次阅读时，不妨就按这种方式来读。我们还要说明：稍后（在 §2.2.3）我们会把导出交换环的术语从「单纯交换环」改为动态交换环。

\subsubsection{导出交换环的其他模型}

我们简要提及文献中常见的「导出交换环」的其他可能模型。

\begin{itemize}[leftmargin=1.6em]
\item \textbf{交换微分分次代数（CDGA）。}这类模型最具体、最便于「显式」写出。例如，用生成元与关系写出 CDGA 相对容易；而对几乎任何单纯交换环，这样的表示都会庞大到写不下来。从我们的角度看，CDGA 的致命缺陷是只在特征 0 中好用，而我们肯定要在整系数下或在特征 $p$ 中工作。不过在特征 0 中，CDGA 完全可以和单纯交换环一样使用。
\item \textbf{$E_\infty$-代数。}粗略地说，$E_\infty$-代数是一个乘法在同伦融贯的意义下交换的环。$E_\infty$-代数的精确定义需要比单纯交换环更多的同伦论。在实践中，它们通常更适合同伦论而非代数几何。确实，$E_\infty$-代数\tn{底本此处作「$E_\infty$ lead to a functional theory」，其中 $E_\infty$ 后疑脱 algebras，functional 疑为 functorial 之误。}能给出对所有特征都适用的函子性理论，在特征 0 时它们也可以和单纯交换环一样使用；但在特征 0 之外，它们的行为不同，而这种不同似乎与我们关系不大。
\end{itemize}

\ysrc{476}
\subsubsection{导出交换环的范畴}

对于模空间理论来说，本质的不是单纯交换环这个概念本身，而是它所在的环境范畴的正确概念。按照「单纯交换环类似于拓扑交换环」的直觉，我们要做的是取拓扑交换环范畴模去同伦等价；非正式地说，我们想把「同伦的」东西视为等价。

为此，我们想与下面这些对象作类比：阿贝尔群、阿贝尔群的复形、阿贝尔群的导出范畴。在这个类比中，交换环对应于阿贝尔群。要「导出」阿贝尔群范畴上的函子，人们考虑「单纯阿贝尔群」，而这类对象的范畴等价于连通链复形的范畴。然而，对许多目的（例如定义导出函子）而言，人们真正关心的并不是阿贝尔群的范畴，而是阿贝尔群的导出范畴。非正式地说，可以这样构造这个范畴：对象是链复形，态射则由链复形之间的态射把拟同构取逆而得到。

单纯交换环与阿贝尔群的链复形相对应。我们真正想要的范畴称为单纯交换环的 $\infty$-范畴，它不是字面意义上的单纯交换环范畴，而是由后者把「弱等价」取逆得到的范畴；弱等价应当看成单纯交换环之间拟同构的类比。

沿用 Clausen 提出、并在 [vS24, §5.1] 中正式解释的术语，我们把上一段中想要的范畴称为动态交换环范畴；在思考该范畴的对象时，我们使用「动态交换环」一词。于是，动态交换环与单纯交换环是同一个对象，但这一术语暗示了所处环境范畴的不同，因而也暗示了这类对象之间（同）构概念的不同。例如，动态交换环的同构对应于单纯交换环之间或许可称为「拟同构」的东西，而不是附录中定义的单纯交换环之间更严格的同构概念。

应当说明，「动态交换环」的概念并不新，文献中传统上就以「单纯交换环」之名处理它。较近时期（至少在本综述涉及的领域中），人们更常用动态这一术语来区分 $\infty$-范畴版本。动态对其他代数结构也普遍适用，例如阿贝尔群：对阿贝尔群作动态，得到的概念与（连通）链复形等价。

\subsection{想象导出概形}

在代数几何中，人们学会从几何上思考交换环。例如，在概形理论的入门教材（如 [Vak24]，其 §4.2 启发了本小节的标题）中，人们学到：环的非约化性——它起初显得「不几何」，因为在域值点层面看不见它——可以在几何上描绘为「无穷小绒毛」（infinitesimal fuzz）。换句话说，人们把 $\Spec R$ 想象为 $\Spec R_{\mathrm{red}}$ 的无穷小加厚。本小节将说明，如何以类似方式把导出概形想象为一种\ysrc{477}经典概形的无穷小加厚。

引导性的口号是：

\begin{quote}
导出概形与经典概形的关系，类似于（经典）概形与约化概形的关系。
\end{quote}

我们先解释形式上的相似之处。一个概形有「底约化」概形，其构造具有函子性，并定义了从约化概形范畴到全体概形范畴的包含函子的右伴随。

\begin{equation}\tag{2.3.1}\label{eq:a6-01-2-3-1}
\begin{tikzcd}[column sep=huge,row sep=small]
\{\text{概形}\} \arrow[rr, bend right=42, "Y\mapsto Y_{\mathrm{red}}"'] & & \{\text{约化概形}\} \arrow[ll, bend right=42, "X\leftarrow X"']
\end{tikzcd}
\end{equation}

在环的层面上，底约化概形的构造局部地对应于把环商去其幂零根；这一构造是遗忘函子的左伴随，后者把约化交换环范畴包含到全体交换环范畴中。

\begin{equation}\tag{2.3.2}\label{eq:a6-01-2-3-2}
\begin{tikzcd}[column sep=huge,row sep=small]
\{\text{交换环}\} \arrow[rr, bend left=42, "R\mapsto R/\mathrm{Nil}"] & & \{\text{约化交换环}\} \arrow[ll, bend left=42, "S\leftarrow S"]
\end{tikzcd}
\end{equation}

现在解释动态交换环的类似图景。任何交换环 $R$ 都可以看作一个动态交换环 $\underline{R}$，直觉上就是「赋予它离散拓扑」。此时 $\pi_0(\underline{R})=R$，而当 $i>0$ 时 $\pi_i(\underline{R})=0$。这给出从交换环范畴到动态交换环范畴的全忠实嵌入，并特别印证了如下观点：动态交换环构成交换环的「扩大」。

另一方面，对任何动态交换环 $\mathcal R$，它的第 0 个同伦群 $\pi_0(\mathcal R)$ 带有自然的交换环结构。这些函子构成一个类似于 (2.3.2) 的伴随\tn{底本此处上方标签的首个符号作斜体 $R$，按上下文疑为动态交换环 $\mathcal R$ 之误。}：

\begin{equation}\tag{2.3.3}\label{eq:a6-01-2-3-3}
\begin{tikzcd}[column sep=huge,row sep=small]
\{\text{动态交换环}\} \arrow[rr, bend left=42, "R\mapsto\pi_0(\mathcal R)"] & & \{\text{交换环}\} \arrow[ll, bend left=42, "\underline{R}\leftarrow R"]
\end{tikzcd}
\end{equation}

这一构造可以粘接起来，于是任何导出概形 $Y$ 都有经典截断 $Y_{\mathrm{cl}}$；其构造具有函子性，并定义了从概形到导出概形的函子的右伴随：

\begin{equation}\tag{2.3.4}\label{eq:a6-01-2-3-4}
\begin{tikzcd}[column sep=huge,row sep=small]
\{\text{导出概形}\} \arrow[rr, bend right=42, "Y\mapsto Y_{\mathrm{cl}}"'] & & \{\text{概形}\} \arrow[ll, bend right=42, "X\leftarrow X"']
\end{tikzcd}
\end{equation}

\ysrc{478}
类似于 (2.3.1)，人们可以把导出概形想象为一个经典概形再加上一些「导出无穷小绒毛」（derived infinitesimal fuzz）。\footnote{实践中的一个差别是：「导出绒毛」可以带有负的（虚拟）维数。这类似于把维数自然延拓到复形上就得到 Euler 示性数，而它可以为负。}

\begin{exam}{2.3.1}\label{exam:a6-01-1}
我们尚未定义导出概形的平展位点，但正如一个概形的平展位点与其底约化概形的平展位点同构，导出概形的平展位点也与它的经典截断的平展位点同构。这进一步支持了导出结构与非约化结构之间的类比。
\end{exam}

此后我们用 $\mathcal A,\mathcal B,\mathcal R$ 这样的花体字母表示单纯交换环，用 $A,B,R$ 这样的正体字母表示经典交换环。根据上面的说明，我们可以把经典交换环看作动态交换环，但仍然偏好用这一约定来强调一个动态交换环来自经典交换环。类似地，我们用 $\mathrm M$ 或 $M$ 这样的正体字母表示经典模空间，用 $\mathcal M$ 这样的花体字母表示导出模空间。

\subsection{单纯交换环的特征}

我们想就如何表达一个动态交换环作一些说明。按照「动态交换环就是同伦意义下的拓扑交换环」这一直觉，自然可以把这个问题与「同伦意义下表达一个拓扑空间」相比较。拓扑空间是一种无定形的对象，但代数拓扑提供了用代数不变量度量它的两种途径：同伦群与（奇异）同调群。我们将解释两个平行的不变量，用来「度量」动态交换环：同伦群与余切复形（cotangent complex）。

\subsubsection{同伦群}

动态交换环 $\mathcal R$ 有一列同伦群 $\pi_0(\mathcal R),\pi_1(\mathcal R),\dots$，它们都是阿贝尔群。不仅如此，$\mathcal R$ 上的乘法使
\[
\pi_*(\mathcal R):=\bigoplus_{i=0}^{\infty}\pi_i(\mathcal R)
\]
具有分次交换环的额外结构。具体构造见 A.5.1。按照「动态交换环就像拓扑交换环」这一直觉类比，你可以把 $\pi_*(\mathcal R)$ 想成底层拓扑空间的同伦群，其环结构由 $\mathcal R$ 的乘法结构诱导。

\begin{exam}{2.4.1}\label{exam:a6-01-2}
若 $\mathcal R$ 来自一个拓扑交换环，则 $\pi_*(\mathcal R)$ 与这个拓扑交换环的同伦群（按通常拓扑意义定义，并以零元素为基点）构成的分次交换环一致。
\end{exam}

对拓扑空间 $X$ 而言，$X$ 的同伦群是「理解」$X$ 本身的好尺度。例如，若 $X$ 与 $X'$ 是「性质良好」的拓扑空间，则态射 $f:X\to X'$ 是同伦等价，当且仅当它对所有 $i$ 都在 $\pi_i$ 上诱导双射——这里要对 $\pi_i$ 作适当解释，以计入连通分支。

\ysrc{479}
类似地，我们把 $\pi_i(\mathcal R)$ 看作「理解」动态交换环 $\mathcal R$ 的好近似。特别地，动态交换环的态射 $f:\mathcal R\to\mathcal R'$ 是同构，当且仅当它对所有 $i$ 都在 $\pi_i$ 上诱导双射。

\subsubsection{余切复形}

设 $f:A\to B$ 是经典交换环的态射。此时有余切复形 $\mathbf L_f$，其定义见 [Ill71,Ill72]，最初是为形变理论的应用而引入的。对我们而言，我们把 $\mathbf L_f$ 看作一个动态 $B$-模。我们在此不给出定义，因为它需要相当长的铺垫；附录 A 对余切复形及其应用作了快速介绍。更一般地，若 $f$ 是概形、导出概形、乃至「性质良好」的叠的态射，则 $f$ 也有余切复形 $\mathbf L_f$。

若 $A=\Z$，我们记 $\mathbf L_B$ 为 $B$ 关于从 $\Z$ 出发的唯一映射的余切复形。更一般地，对概形 $X$（或导出概形、或「性质良好」的叠），我们记 $\mathbf L_X$ 为唯一映射 $X\to\Spec\Z$ 的余切复形。

\begin{exam}{2.4.2}\label{exam:a6-01-3}
设 $f:A\to B$ 是交换环的态射，$\wt A\to A$ 是理想为 $I$ 的平方零加厚（即 $I^2=0$）。令 $J=I\otimes_A B$。我们考虑 $B$ 的 $\wt A$-代数形变，即平坦的 $\wt A$-代数 $\wt B$，使得 $\wt B\otimes_{\wt A}\wt A/I\isom B$。§A.15 说明了这一形变论问题由 $\mathbf L_f$「控制」，具体含义是：
\begin{itemize}[leftmargin=1.6em]
\item 存在类 $\mathrm{obs}(I)\in\Ext^2_B(\mathbf L_f,J)$，它消失当且仅当这样的形变 $\wt B$ 存在。
\item 若 $\mathrm{obs}(I)=0$，则形变 $\wt B$ 的集合具有自然的 $\Ext^1_B(\mathbf L_f,J)$-挠子（torsor）结构。
\item 任何形变 $\wt B$ 的自同构群自然同构于 $\Ext^0_B(\mathbf L_f,J)\cong\Hom_B(\Omega_f,J)$。
\end{itemize}
\end{exam}

\begin{exam}{2.4.3}\label{exam:a6-01-4}
若 $f$ 光滑，则 $\mathbf L_f$ 的上同调集中在 0 次，即由一个 $B$-模表示，而这个模不是别的，正是 Kähler 微分 $\Omega_{B/A}$。
\end{exam}

事实上，余切复形的构造用到了单纯交换环，这也是该理论最早的应用之一。粗略地说，想法是：当 $B$ 在 $A$ 上不光滑时，用一个光滑的单纯 $A$-代数来「消解」$B$，再取这个消解的 Kähler 微分。这令人想起同调代数中定义导出函子的方式：用具有良好性质的某类链复形「消解」输入。链复形只对阿贝尔范畴的对象有意义，而单纯对象的意义要广泛得多。

\begin{exam}{2.4.4}\label{exam:a6-01-5}
设给定一列映射
\begin{equation*}
\begin{tikzcd}[column sep=large]
A \arrow[r, "f"] \arrow[rr, bend left=32, "h"] & B \arrow[r, "g"] & C
\end{tikzcd}
\end{equation*}
其中 $h=g\circ f$，则得到正合列
\[
\Omega_f\otimes_B C\to\Omega_h\to\Omega_g\to 0.
\]
\end{exam}

\ysrc{480}
这一正合列可以延拓为余切复形的正合三角（在 $C$-模范畴的导出范畴中）
\[
\mathbf L_f\otimes^{\mathbf L}_B C\to\mathbf L_h\to\mathbf L_g\xrightarrow{+1}
\]
它特别给出上同调中的长正合列
\[
\dots\to H^{-1}(\mathbf L_h)\to H^{-1}(\mathbf L_g)\to\Omega_f\otimes_B C\to\Omega_h\to\Omega_g\to 0.
\]

按定义，动态交换环的态射是同构，当且仅当它在同伦群上诱导同构。相比之下，余切复形究竟携带「多少」信息，先验地并不清楚。结果表明它携带导出结构的「全部」信息。下面的 Lurie 引理把这一点讲清楚了：

\begin{lem}{2.4.5}（[Lur04, Corollary 3.2.17]）\label{lem:a6-01-1}
设单纯交换环的态射 $f:\mathcal A\to\mathcal B$ 在 $\pi_0$ 上诱导同构，并诱导同构
\[
\mathbf L_{\mathcal A}\otimes_{\mathcal A}\mathcal B\isom\mathbf L_{\mathcal B}.
\]
则 $f$ 是同构。
\end{lem}

我们不给出引理 2.4.5 的形式论证，但会勾勒其直觉。从函子观点看，要证明的是：$f$ 在从 $\mathcal A$ 出发与从 $\mathcal B$ 出发到任意单纯交换环的态射之间诱导等价。如 §2.3 所述，任意单纯交换环都可以看作其经典截断的「导出幂零加厚」。由于到经典环的态射由 $\pi_0$ 控制，从 $\mathcal A$ 出发与从 $\mathcal B$ 出发到任何经典环的态射被 $f$ 等同起来。然后，由例 2.4.2（的推广），到经典环的导出幂零加厚的态射由各自的余切复形控制，而这些余切复形同样被 $f$ 等同起来。

\begin{rem}{2.4.6}\label{rem:a6-01-1}
我们把余切复形看作单纯交换环的「同调理论」。从这个角度看，引理 2.4.5 类似于 Hurewicz 定理：该定理说，对「性质良好」的空间，一个映射在各阶同伦群上诱导双射，与在同调群上诱导双射是等价的。
\end{rem}

\subsubsection{切复形}

设 $f:\mathcal A\to\mathcal B$ 是动态交换环的态射，其余切复形 $\mathbf L_f$ 作为动态 $B$-模是完美的。这意味着 $\mathbf L_f$ 可以局部地由自由 $B$-模的有限复形表示。于是我们定义 $f$ 的切复形（tangent complex）为
\[
\mathbf T_f=\mathcal{RH}om_{\mathcal B}(\mathbf L_f,\mathcal B),
\]
即余切复形的（导出）对偶。

这一构造可以整体化：若 $f:\mathcal Y\to\mathcal X$ 是（导出）概形（或性质良好的叠，使得 $\mathbf L_f$ 存在）的态射，且 $\mathbf L_f$ 在 $\mathcal O_{\mathcal Y}$ 上完美，则定义
\[
\mathbf T_f:=\mathcal{RH}om_{\mathcal Y}(\mathbf L_f,\mathcal O_{\mathcal Y}).
\]

\subsection{导出概形的例子}

仿射导出概形反等价于动态交换环。对动态交换环 $\mathcal A$，我们记 $\Spec\mathcal A$ 为相应的仿射导出概形。一般的导出概形由仿射导出概形粘接而成，类似于概形由仿射概形粘接而成。下面给出实践中导致导出概形的一些自然构造的例子。

\subsubsection{导出纤维积}

即使从纯经典概形出发，导出概形也会通过导出纤维积运算自然地出现。

\ysrc{481}
这一运算的经典化身是概形态射 $X\to Z$ 与 $Y\to Z$ 的通常纤维积。所得的纤维积 $X\times_Z Y$ 在带有到 $X$ 与 $Y$ 的映射、且到 $Z$ 的复合映射彼此一致的概形中，由万有性质刻画。它在仿射局部由交换环的张量积运算构造。

类似地，导出纤维积——我们有时也称它为同伦纤维积 $X\times^h_Z Y$——由类似的万有性质刻画。它在仿射局部由动态交换环上的导出张量积运算构造（见下文）。更一般地，同样的考虑可以构造导出概形的导出纤维积。

给定动态交换环的态射 $C\to A$ 与 $C\to B$，导出张量积 $A\otimes^{\mathbf L}_C B$ 可以由其万有性质刻画：它是动态 $C$-代数范畴中 $A$ 与 $B$ 的余积。用单纯交换环可以给出一个显式模型：选取 $A$ 与 $B$ 作为 $C$-代数的「消解」，详见 §A.10。特别地，这一显式模型表明 $A\otimes^{\mathbf L}_C B$ 的同伦群就是 Tor 群：
\begin{equation}\tag{2.5.1}\label{eq:a6-01-2-5-1}
\pi_i(A\otimes^{\mathbf L}_C B)\cong\Tor^C_i(A,B);
\end{equation}
不过 (2.5.1) 并未揭示 $\pi_*(A\otimes^{\mathbf L}_C B)$ 上的乘法结构，所以它实际上只编码了 $A\otimes^{\mathbf L}_C B$ 底层动态阿贝尔群的信息。由 (2.5.1) 可知，经典概形的导出纤维积 $X\times^h_Z Y$ 是经典的，当且仅当 $\mathcal O_X$ 与 $\mathcal O_Y$ 在 $\mathcal O_Z$ 上 Tor-无关。特别地，若 $X\to Z$ 或 $Y\to Z$ 是平坦的，则此条件成立。

\begin{exam}{2.5.1}\label{exam:a6-01-6}
导出纤维积与（余）切复形的构造相容得很好，这实际上是导出代数几何的一个技术优势。更一般地，有一个口号：「切复形的构造与同伦极限交换」。特别地，对导出概形 $\mathcal X\to\mathcal Z\leftarrow\mathcal Y$，有
\[
\mathbf T_{\mathcal X\times_{\mathcal Z}\mathcal Y}\cong\mathrm{Fib}\big(\mathbf T_{\mathcal X}|_{\mathcal X\times_{\mathcal Z}\mathcal Y}\oplus\mathbf T_{\mathcal Y}|_{\mathcal X\times_{\mathcal Z}\mathcal Y}\to\mathbf T_{\mathcal Z}|_{\mathcal X\times_{\mathcal Z}\mathcal Y}\big),
\]
其中纤维在 $\mathcal X\times_{\mathcal Z}\mathcal Y$ 上的拟凝聚层导出范畴中取（这一记号即指导出纤维积，因为对导出概形而言并不存在「朴素纤维积」）。
\end{exam}

\subsubsection{导出向量丛}

设 $X$ 是域 $k$ 上的本征簇，$\mathcal E$ 是 $X$ 上的局部自由凝聚层。考虑把 $S\in\mathrm{Sch}_{/k}$ 送到 $\mathcal E$ 在 $X_S$ 上整体截面的空间的函子；它由
\[
\underline{H^0(X,\mathcal E)}=H^0(X,\mathcal E)\otimes_k\mathbf A^1_k
\]
表示，即把 $k$-向量空间 $H^0(X,\mathcal E)$ 看作 $k$ 上的仿射空间。它也可以描述为 $\Spec(\Sym H^0(X,\mathcal E)^*)$，其中 $H^0(X,\mathcal E)^*$ 是 $H^0(X,\mathcal E)$ 在 $k$ 上的对偶。

同调代数教给我们的一个视角是：把 $H^0(X,\mathcal E)$ 看作复形 $\mathrm{R}\Gamma(X,\mathcal E)$ 的第零个上同调群，而后者更高次的上同调群是 $H^i(X,\mathcal E)$（$i\ge1$）。类似地，有导出仿射概形
\[
\underline{\mathrm{R}\Gamma}(X,\mathcal E)=\text{“}\operatorname{Spec}\big(\operatorname{Sym}\mathrm{R}\Gamma(X,\mathcal E)^*\big)\text{”},
\]
其经典截断为 $\underline{H^0(X,\mathcal E)}$。它表示如下函子：非正式地说，它把导出概形 $S$ 送到 $\mathrm{R}\Gamma(X_S,\mathcal E)$。我们在这里看到的是：导出代数几何让我们能够构造一些导出概形，它们\ysrc{482}延拓整体截取的经典模空间，正如导出函子延拓经典的整体截取。

现在，我们在上面的公式中加了引号，因为我们还没有说明动态 $k$-模的对称幂是什么意思。它可以刻画为遗忘函子的左伴随，该遗忘函子从动态交换 $k$-代数映到动态 $k$-模。

\begin{rem}{2.5.2}（导出向量丛）\label{rem:a6-01-2}
更一般地，可以做「相对于基」的类似构造：给定（导出）叠 $S$ 上由拟凝聚层构成的完美复形 $\mathcal K$，可以构造作为 $S$ 上导出叠的「全空间」$\operatorname{Tot}_S(\mathcal K)$。这类构造在 [FYZ23] 中被称为导出向量丛，因为它们推广了向量丛：向量丛正是 $\mathcal K$ 为局部自由凝聚层（看作集中在 0 次的完美复形）时出现的全空间。
\end{rem}

\begin{exam}{2.5.3}\label{exam:a6-01-7}
设 $S=\mathrm{Bun}_n$ 是光滑射影曲线 $C$ 上秩 $n$ 向量丛的模叠。考虑空间 $\mathrm M\to S$，其 $R$-点是 $C$ 上秩 $n$ 向量丛 $\mathcal F$ 连同 $\mathcal F$ 的一个整体截面构成的群胚。虽然 $\mathrm M\to S$ 的纤维是向量空间，态射 $\mathrm M\to S$ 甚至不是平坦的，因为纤维的维数不同（即使在基的连通分支上限制也是如此）。这是因为 $\mathcal F\mapsto H^0(\mathcal F)$ 的行为「不连续」，例如 $H^0(\mathcal F)$ 的维数随 $\mathcal F$ 不连续地变化。

然而，完美复形 $\mathrm R\Gamma(\mathcal F)$ 的行为「连续」，例如 $\mathrm R\Gamma(\mathcal F)$ 的 Euler 示性数随 $\mathcal F$ 连续变化。这预示了如下事实：我们可以把 $R\Gamma(\mathcal F)$ 组装成 $S$ 上的完美复形，其全空间 $\mathcal M\to S$ 的经典截断是 $\mathrm M\to S$，而它在 $\mathcal F$ 上的导出纤维是上述意义下的导出概形 $R\Gamma(\mathcal F)$。这个 $\mathcal M$ 是 §3 与 §4 将要讨论的「导出 Hitchin 叠」的一个变体。
\end{exam}

\subsubsection{导出模空间}

导出概形（或叠）最有趣的例子大概是导出模空间。§3 整节都用来讨论这样的例子，所以这里只给一个简短的预告。这里的想法是：先从一个经典模空间出发——按定义，它是定义在交换环上、满足适当下降条件的函子。然后找到一种方式，尽可能重新表述模问题，使它在动态交换环上也有意义，并且一般输出动态集（anima，即动态化的集合）而非集合。表示这一函子的对象称为导出模空间。稍后我们会看到把模空间提升为导出模空间的一些好处。

\subsection{拟光滑性}

本小节定义导出概形态射的一个性质，称为拟光滑性（quasismoothness）；它在导出代数几何中起着非常重要的作用。它是局部完全交态射的导出推广。作为动机，我们回顾一下局部完全交的含义：

\begin{defn}{2.6.1}\label{defn:a6-01-1}
设 $A$ 是诺特环，$f:A\to B$ 是有限型态射。若 $f$ 在 $A$ 上 Zariski 局部地是如下映射的复合：先作形如 $A\to A[x_1,\dots,x_n]$ 的映射，再商去一个正则序列，则称 $f$ 为局部完全交（local complete intersection，LCI）。
\end{defn}

局部完全交可以用余切复形给出一个干净利落的刻画。
